% GENERATED FILE. Edit canonical lessons, then run npm run generate:books.
% canonical-source-hash: fnv1a64:7355416f19be2921
% canonical-lessons: SA-C321-abhisekah, SA-C321-purohitah, SA-C321-mantrah, SA-C321-abhisincati, SA-C321-uccarayati

\chapter{A Coronation}
\label{ch:sa-a-coronation}
\hlchaptermodality{321}

\begin{chapteropening}
\textbf{By the end of this chapter:} \emph{I can describe how a king was crowned.}

Complete the last lesson of chapter 321: I can describe how a king was crowned.
\end{chapteropening}

\section[\texorpdfstring{\sk{अभिषेकः} (abhiṣekaḥ)}{अभिषेकः (abhiṣekaḥ)}]{\sk{अभिषेकः} (abhiṣekaḥ) — a coronation; an anointing}

\label{lesson:SA-C321-abhisekah}



\begin{quote}
Before the new one: say the Sanskrit for drives out, then the Sanskrit for freedom.
\end{quote}

\subsection*{You'll want to know: \sk{अभिषेकः}}
\textbf{\sk{अभिषेकः}} — \emph{abhiṣekaḥ} — \textquotedblleft{}a coronation; an anointing\textquotedblright{}.

\begin{quote}
\sk{अद्य} \sk{प्रासादे} \sk{अभिषेकः} \sk{भवति।}
\end{quote}

\begin{quote}
\emph{adya prāsāde abhiṣekaḥ bhavati.}
\end{quote}

\begin{quote}
Today there is a coronation in the palace.
\end{quote}

A masculine noun.

\begin{cousinweb}
\textbf{\sk{अभिषेकः}} (\emph{abhiṣekaḥ}) is built from \emph{abhi} and the root √sic, \textquotedblleft{}to sprinkle\textquotedblright{}; MW gives \emph{abhiṣeka} as anointing or consecrating by sprinkling water, and the inauguration of a king (MW).
\end{cousinweb}

\begin{grammarlens}[title={What you've built}]
One more word for telling what happened.
\end{grammarlens}

\subsection*{Your turn}
\begin{itemize}
\raggedright
  \item \emph{Say it:} \emph{abhiṣekaḥ}
  \item \emph{Say it:} \emph{abhiṣekaḥ}, once more
  \item \emph{Say it:} \emph{adya prāsāde abhiṣekaḥ bhavati.} — Today there is a coronation in the palace.
  \item \emph{Recall:} say the Sanskrit for drives out, then the Sanskrit for freedom, then say \emph{abhiṣekaḥ} again
\end{itemize}

\begin{tcolorbox}[breakable,colback=teal!4,colframe=teal!35!black,title={Before you move on}]
What does \sk{अभिषेकः} mean? (\textquotedblleft{}A coronation; an anointing\textquotedblright{}.) Say it once more.
\end{tcolorbox}
\section[\texorpdfstring{\sk{पुरोहितः} (purohitaḥ)}{पुरोहितः (purohitaḥ)}]{\sk{पुरोहितः} (purohitaḥ) — a family priest}

\label{lesson:SA-C321-purohitah}



\begin{quote}
Before the new one: say the Sanskrit for freedom, then the Sanskrit for a coronation.
\end{quote}

\subsection*{You'll want to know: \sk{पुरोहितः}}
\textbf{\sk{पुरोहितः}} — \emph{purohitaḥ} — \textquotedblleft{}a family priest\textquotedblright{}.

\begin{quote}
\sk{पुरोहितः} \sk{अभिषेकं} \sk{करोति।}
\end{quote}

\begin{quote}
\emph{purohitaḥ abhiṣekaṃ karoti.}
\end{quote}

\begin{quote}
The priest performs the coronation.
\end{quote}

A masculine noun.

\begin{cousinweb}
\textbf{\sk{पुरोहितः}} (\emph{purohitaḥ}) is built from \emph{puras}, \textquotedblleft{}in front\textquotedblright{}, and \emph{hita}, \textquotedblleft{}placed\textquotedblright{}, from the root √dhā; MW gives \emph{purohita} as placed in front, appointed, and especially a family priest (MW).
\end{cousinweb}

\begin{grammarlens}[title={What you've built}]
One more word for telling what happened.
\end{grammarlens}

\subsection*{Your turn}
\begin{itemize}
\raggedright
  \item \emph{Say it:} \emph{purohitaḥ}
  \item \emph{Say it:} \emph{purohitaḥ}, once more
  \item \emph{Say it:} \emph{purohitaḥ abhiṣekaṃ karoti.} — The priest performs the coronation.
  \item \emph{Recall:} say the Sanskrit for freedom, then the Sanskrit for a coronation, then say \emph{purohitaḥ} again
\end{itemize}

\begin{tcolorbox}[breakable,colback=teal!4,colframe=teal!35!black,title={Before you move on}]
What does \sk{पुरोहितः} mean? (\textquotedblleft{}A family priest\textquotedblright{}.) Say it once more.
\end{tcolorbox}
\section[\texorpdfstring{\sk{मन्त्रः} (mantraḥ)}{मन्त्रः (mantraḥ)}]{\sk{मन्त्रः} (mantraḥ) — a sacred verse; a mantra}

\label{lesson:SA-C321-mantrah}



\begin{quote}
Before the new one: say the Sanskrit for a coronation, then the Sanskrit for a family priest.
\end{quote}

\subsection*{You'll want to know: \sk{मन्त्रः}}
\textbf{\sk{मन्त्रः}} — \emph{mantraḥ} — \textquotedblleft{}a sacred verse; a mantra\textquotedblright{}.

\begin{quote}
\sk{पुरोहितः} \sk{मन्त्रं} \sk{पठति।}
\end{quote}

\begin{quote}
\emph{purohitaḥ mantraṃ paṭhati.}
\end{quote}

\begin{quote}
The priest recites a sacred verse.
\end{quote}

A masculine noun.

\begin{cousinweb}
\textbf{\sk{मन्त्रः}} (\emph{mantraḥ}) comes from the root √man, \textquotedblleft{}to think\textquotedblright{}: MW explains \emph{mantra} as an instrument of thought, and gives it as speech, a sacred text, a prayer or hymn (MW).
\end{cousinweb}

\begin{grammarlens}[title={What you've built}]
One more word for telling what happened.
\end{grammarlens}

\subsection*{Your turn}
\begin{itemize}
\raggedright
  \item \emph{Say it:} \emph{mantraḥ}
  \item \emph{Say it:} \emph{mantraḥ}, once more
  \item \emph{Say it:} \emph{purohitaḥ mantraṃ paṭhati.} — The priest recites a sacred verse.
  \item \emph{Recall:} say the Sanskrit for a coronation, then the Sanskrit for a family priest, then say \emph{mantraḥ} again
\end{itemize}

\begin{tcolorbox}[breakable,colback=teal!4,colframe=teal!35!black,title={Before you move on}]
What does \sk{मन्त्रः} mean? (\textquotedblleft{}A sacred verse; a mantra\textquotedblright{}.) Say it once more.
\end{tcolorbox}
\section[\texorpdfstring{\sk{अभिषिञ्चति} (abhiṣiñcati)}{अभिषिञ्चति (abhiṣiñcati)}]{\sk{अभिषिञ्चति} (abhiṣiñcati) — anoints; consecrates}

\label{lesson:SA-C321-abhisincati}



\begin{quote}
Before the new one: say the Sanskrit for a family priest, then the Sanskrit for a sacred verse.
\end{quote}

\subsection*{You'll want to know: \sk{अभिषिञ्चति}}
\textbf{\sk{अभिषिञ्चति}} — \emph{abhiṣiñcati} — \textquotedblleft{}anoints; consecrates\textquotedblright{}.

\begin{quote}
\sk{पुरोहितः} \sk{तं} \sk{जलेन} \sk{अभिषिञ्चति।}
\end{quote}

\begin{quote}
\emph{purohitaḥ taṃ jalena abhiṣiñcati.}
\end{quote}

\begin{quote}
The priest anoints him with water.
\end{quote}

Past (he …): \sk{सः} \sk{तं} \sk{जलेन} \sk{अभिषिक्तवान्}, \emph{saḥ taṃ jalena abhiṣiktavān}.

\begin{cousinweb}
\textbf{\sk{अभिषिञ्चति}} (\emph{abhiṣiñcati}) is built from \emph{abhi} and the root √sic, \textquotedblleft{}to sprinkle\textquotedblright{}; MW gives \emph{abhi-ṣic} as to sprinkle and to consecrate, anoint, appoint by consecration (MW).
\end{cousinweb}

\begin{grammarlens}[title={What you've built}]
One more word for telling what happened.
\end{grammarlens}

\subsection*{Your turn}
\begin{itemize}
\raggedright
  \item \emph{Say it:} \emph{abhiṣiñcati}
  \item \emph{Say it:} \emph{abhiṣiñcati}, once more
  \item \emph{Say it:} \emph{purohitaḥ taṃ jalena abhiṣiñcati.} — The priest anoints him with water.
  \item \emph{Recall:} say the Sanskrit for a family priest, then the Sanskrit for a sacred verse, then say \emph{abhiṣiñcati} again
\end{itemize}

\begin{tcolorbox}[breakable,colback=teal!4,colframe=teal!35!black,title={Before you move on}]
What does \sk{अभिषिञ्चति} mean? (\textquotedblleft{}Anoints; consecrates\textquotedblright{}.) Say it once more.
\end{tcolorbox}
\section[\texorpdfstring{\sk{उच्चारयति} (uccārayati)}{उच्चारयति (uccārayati)}]{\sk{उच्चारयति} (uccārayati) — pronounces; utters}

\label{lesson:SA-C321-uccarayati}



\begin{quote}
Before the new one: say the Sanskrit for a sacred verse, then the Sanskrit for anoints.
\end{quote}

\subsection*{You'll want to know: \sk{उच्चारयति}}
\textbf{\sk{उच्चारयति}} — \emph{uccārayati} — \textquotedblleft{}pronounces; utters\textquotedblright{}.

\begin{quote}
\sk{बालः} \sk{श्लोकं} \sk{स्पष्टम्} \sk{उच्चारयति।}
\end{quote}

\begin{quote}
\emph{bālaḥ ślokaṃ spaṣṭam uccārayati.}
\end{quote}

\begin{quote}
The boy pronounces the verse clearly.
\end{quote}

Past (he …): \sk{सः} \sk{मन्त्रम्} \sk{उच्चारितवान्}, \emph{saḥ mantram uccāritavān}.

\begin{cousinweb}
\textbf{\sk{उच्चारयति}} (\emph{uccārayati}) is built from \emph{ud}, \textquotedblleft{}up, out\textquotedblright{}, and the causative of the root √car, \textquotedblleft{}to move\textquotedblright{}; MW gives \emph{uc-cārayati} as to cause to go forth, to utter, pronounce (MW).
\end{cousinweb}

\begin{grammarlens}[title={What you've built}]
One more word for telling what happened.
\end{grammarlens}

\subsection*{Your turn}
\begin{itemize}
\raggedright
  \item \emph{Say it:} \emph{uccārayati}
  \item \emph{Say it:} \emph{uccārayati}, once more
  \item \emph{Say it:} \emph{bālaḥ ślokaṃ spaṣṭam uccārayati.} — The boy pronounces the verse clearly.
  \item \emph{Recall:} say the Sanskrit for a sacred verse, then the Sanskrit for anoints, then say \emph{uccārayati} again
\end{itemize}

\begin{tcolorbox}[breakable,colback=teal!4,colframe=teal!35!black,title={Before you move on}]
What does \sk{उच्चारयति} mean? (\textquotedblleft{}Pronounces; utters\textquotedblright{}.) Say it once more.
\end{tcolorbox}
